\documentclass[12pt]{elsarticle}

\usepackage[margin=1in]{geometry}
\usepackage[hidelinks]{hyperref}
\usepackage[utf8]{inputenc}
\usepackage[english]{babel}

\begin{document}
\thispagestyle{empty}

\begin{center}
    {\textbf{Resolving Coordination Frictions in Green Labor Transitions:}}\\[0.4em]
    { \textbf{Minimizing Unemployment, Costs, and Welfare Distortions}}\\[2em]

    {\normalsize
    Shisham Adhikari\footnote{Corresponding author. Email: \texttt{shadhikari@ucdavis.edu}}\\[0.5em]
    Department of Economics, University of California, Davis\\
    One Shields Avenue, Davis, CA 95616, USA}
\end{center}

\vspace{1.5em}

\noindent\textbf{Abstract:}
Ambitious decarbonization requires not only technology and capital, but also rapid reallocation of workers into green jobs. This reallocation faces a two-sided coordination failure: workers hesitate to pay entry costs without job prospects, while firms underinvest in vacancies without a ready talent pool. This friction compounds the environmental externality and stalls the transition when both speed and scale are essential. I develop a two-sector Diamond–Mortensen–Pissarides model with costly worker entry, firm vacancy creation, and balanced-budget instruments to compare worker-side support, firm-side support, and coordinated dual policy. Calibrated to a U.S. energy transition from 2\% to 14\% green employment by 2030 (24 million jobs), I find three results. First, coordinated policy dominates one-sided alternatives, reaching the target with 39 basis points lower unemployment, 5 basis points of GDP lower fiscal cost, and 23 basis points higher welfare. Second, optimal policy is dynamic: total support is hump-shaped, and the instrument mix rotates worker–firm–worker as the binding constraint shifts along the transition path. Third, coordinated policy raises aggregate welfare once the environmental benefit exceeds a consumption-equivalent gain of 0.87\%. The main lesson is that climate policy design, not just scale, matters for socially inclusive and fiscally sustainable transitions.

\vspace{1em}

\noindent\textbf{Keywords:}
Climate change; Green jobs; Search and matching frictions; Climate policy design; Just transition

\end{document}
